\documentclass[11pt]{article}

\usepackage[margin=1in]{geometry}
\usepackage{graphicx}
\usepackage{booktabs}
\usepackage{amsmath}
\usepackage{amssymb}
\usepackage{xcolor}
\usepackage[normalem]{ulem}
\usepackage[colorlinks=true,linkcolor=blue,urlcolor=blue]{hyperref}
\usepackage{enumitem}

% Figures are included with explicit fig/<name> paths (self-contained;
% regenerate originals in ../verification/).

\newcommand{\fpl}{\textsc{fpl}}
\newcommand{\mppi}{\textsc{mppi}}
\newcommand{\degree}{\ensuremath{^\circ}}   % works in text and math
\newcommand{\claim}[1]{\textbf{\textcolor{blue!60!black}{#1}}}

\title{\vspace{-2em}Should we publish the \fpl$\times$\mppi{} results?\\
\large A memo, with the evidence laid out and the case argued honestly}
\author{}
\date{\today}

\begin{document}
\maketitle
\vspace{-2.5em}

\begin{abstract}
\noindent This is a memo to myself (and you) making the case, for and against, that the
current results are worth writing up. Short version: \emph{yes}, but the paper should be
built around the \textbf{model-mismatch robustness} result, not the raw performance numbers.
The performance win alone ($+$34--93\% speed/rotation at matched safety across four tasks,
30 seeds) is real but invites the reviewer reflex ``\,so retune your linear weights.\,'' The robustness
result closes that door: a \emph{single fixed} \fpl{} objective dominates the entire
linear-weight family across an aggression-punishing model change with \emph{no retuning}, on
\textbf{four distinct robots (a hopper, a planar biped, a 12-DoF whole-body quadruped, and a
16-DoF dexterous hand)} and \textbf{two mismatch axes (ground traction and contact grip)}
---and you cannot retune a hand-weighted cost for a model error you do not know about. The
dexterous-hand task (in-hand cube reorientation) is \emph{non-locomotion}, so the phenomenon is
not a property of legged gaits. That is a structural, un-tunable property, which is the kind of
thing that gets published. We further show the advantage is a property of the \textbf{objective},
not the search: it holds \emph{within} every sampler we tried (\mppi, \mppi-CMA, CEM, colored
noise), no sampler beats plain \mppi, and a \emph{single} \fpl{} spec transfers across morphologies
with \textbf{no tuning}, while no fixed linear weight does. The honest negative results (a smarter
sampler, a smarter weighting, and a hard safety shield all fail to beat the simple design) make the
story more credible, not less.
\end{abstract}

\section{What we actually claim}

We study a single design choice inside sampling-based model-predictive control (\mppi):
how a vector of objectives is collapsed into the scalar cost that weights the sampled
rollouts. The standard choice is a \textbf{hand-weighted linear sum}
$S = \sum_i w_i c_i$. We replace it with \textbf{Fulfillment-Priority Logic} (\fpl):
each objective is mapped to a fulfillment $f_i\in[0,1]$ and composed with a
power mean $\mu_p(f)=\big(\tfrac1n\sum_i f_i^{\,p}\big)^{1/p}$ using $p<0$
(a soft-min / conjunction that weights the least-satisfied objective most).
\emph{Everything else in the controller is held identical}; only the scalarization changes.
Because $\mu_p(f;p{=}1)=\sum_i w_i f_i$ is exactly the linear sum, the linear family is a
strict special case, and ``the linear-weight family'' is a one-parameter sweep we compare
against \emph{at its best}, not its worst.

\medskip
\noindent Three empirical claims, in increasing order of how much I think they matter:

\begin{enumerate}[leftmargin=1.4em,itemsep=2pt]
\item \claim{Dynamic-competition dominance.} When objectives compete \emph{dynamically}
      (going faster unavoidably risks a fall), one fixed \fpl{} spec Pareto-dominates the
      entire linear-weight family on the (speed, survival) plane, with no per-operating-point
      retuning. Shown on Hopper, Walker2d, a 12-DoF quadruped, and a 16-DoF LEAP hand
      (in-hand cube reorientation---a \emph{non-locomotion} task, so the effect is not about gaits).
\item \claim{An honest boundary.} That dominance \emph{disappears} when a safe conservative
      operating point exists (quasi-static tasks) or the task is unachievable for the
      planner. We show two such nulls. This is a scientific boundary, not a failure.
\item \claim{Un-tunable robustness (the headline).} Under model mismatch that punishes
      aggression (traction loss on the legged robots, contact-grip loss on the hand), the
      single fixed \fpl{} spec keeps the highest \emph{productive} performance
      (progress $\times$ survival) at every mismatch level on all four robots, while every fixed
      linear weight either was already unsafe or degrades---and no retuning is possible because
      the mismatch is unknown at design time.
\end{enumerate}

\section{The evidence}

\paragraph{Fairness protocol.} Identical sampler, fulfillment atoms, sample budget, planning
horizon, and temporal aggregation for \emph{every} configuration. The only variable is the
composition: $p{=}1$ with swept weights (the linear family) versus $p{<}0$ uniform (\fpl).
We always compare against the \emph{best} linear weight, not the aggressive one that trivially
face-plants. On the mismatch study, the controller plans on the nominal model and executes on
a perturbed ``true'' simulator; \emph{no controller is retuned across mismatch levels.}

\paragraph{1. Dynamic-competition dominance (Fig.~\ref{fig:pareto}).}
On all four robots a single \fpl{} point sits above-and-right of the whole linear frontier at
every commanded speed (or, for the hand, commanded rotation). Headline numbers, fastest linear
at $\le$ \fpl's fall (drop) rate:

\begin{center}\small
\begin{tabular}{lccccc}
\toprule
 & Hopper (tv$=3$) & Walker (tv$=4$) & Walker (tv$=6$) & Quadruped (tv$=3$) & Hand (roll$=103\degree$) \\
\midrule
\fpl{} (one fixed spec) & 0.975\,@\,33\% & 1.16\,@\,0\% & 1.20\,@\,0\% & 1.21\,@\,0\% & 51.6\degree\,@\,0\% \\
best equally-safe linear & 0.539 & 0.77 & 0.90 & 0.77 & 26.8\degree \\
\fpl{} advantage & \textbf{+81\%} & \textbf{+50\%} & \textbf{+34\%} & \textbf{+57\%} & \textbf{+93\%} \\
\bottomrule
\end{tabular}\\[2pt]
{\footnotesize achieved speed (m/s) or rotation ($\degree$ toward goal) @ fall/drop rate;
``advantage'' vs.\ the fastest linear weight at $\le$ \fpl's fall/drop rate. The hand column is
the fastest linear weight within $10\%$ drop; the only \emph{strictly} $0\%$-drop linear weight
is the near-static crawler ($\sim$10$\degree$).}
\end{center}

\begin{figure}[t]
\centering
\includegraphics[width=\linewidth]{fig/hopper_pareto.png}\\[3pt]
\includegraphics[width=\linewidth]{fig/walker_pareto.png}\\[3pt]
\includegraphics[width=\linewidth]{fig/quadruped_pareto.png}\\[3pt]
\includegraphics[width=\linewidth]{fig/cube_pareto.png}
\caption{\textbf{One fixed \fpl{} spec dominates the linear-weight family under dynamic
competition, no retuning.} Top to bottom: Hopper; Walker2d (running regime); Barkour
quadruped (12 DoF whole-body); LEAP hand (16 DoF, in-hand cube reorientation---the
\emph{non-locomotion} task). At every commanded speed / rotation the \fpl{} star is faster
\emph{and} safer than any comparably-safe linear weight; the linear family traces a
progress$\leftrightarrow$survival tradeoff that \fpl{} sits above, on all four robots.}
\label{fig:pareto}
\end{figure}

\paragraph{2. The honest boundary (nulls).} On a quasi-static humanoid lean-and-hold task,
a \emph{conservative} linear weight matches or beats \fpl{} at every difficulty, because a safe
low-risk operating point exists. On humanoid walking-from-stand, the planner's short horizon
cannot walk it and \emph{both} controllers collapse. These delimit the claim: ``objectives
compete'' is not sufficient---the competition must be dynamic and unavoidable. (An early
version of this project over-claimed a humanoid win from an uprightness-only metric; rendering
the video exposed a seated collapse. We now co-check height and inspect video before claiming.
That discipline is worth a paragraph in the paper: it is why the remaining claims are trustworthy.)

\paragraph{3. Un-tunable robustness (Fig.~\ref{fig:robust}).} The controller plans on the
nominal model; the true ground friction is reduced (traction loss---ice / wet / wear, the
canonical sim-to-real gap). Productive speed (speed $\times$ survival), no retuning:

\begin{center}\small
\begin{tabular}{lcc}
\toprule
robot & \fpl{} productive speed (range) & best fixed linear (range) \\
\midrule
Hopper       & \textbf{0.66--0.93} & 0.34--0.77 \\
Walker2d     & \textbf{1.09--1.18} & 0.81--0.98 \\
Quadruped    & \textbf{1.09--1.15} & 0.78--1.01 \\
\bottomrule
\end{tabular}\\[2pt]
{\footnotesize Productive speed (m/s $\times$ survival) held across the traction sweep, 30
seeds; \fpl{} is highest at every friction level on all three legged robots (the hand, a
different metric and mismatch axis, is in paragraph~4).}
\end{center}

\noindent The fixed \fpl{} spec is best at \emph{every} mismatch level on all three legged
robots, and its margin on the hopper \emph{grows} with mismatch (+49\% at nominal $\to$ +103\%
at the slipperiest). No fixed linear weight matches it: the
aggressive weights' survival collapses, the safe weight crawls, and on the hopper the
\emph{best} linear weight even changes with the (unknown) friction. On the quadruped the
picture is the same---\fpl{} holds 100\% survival at top-or-tied productive speed while the
aggressive linear weights invert onto the robot's back (Fig.~\ref{fig:robust_quad}).

\begin{figure}[t]
\centering
\includegraphics[width=0.92\linewidth]{fig/mismatch_robustness.png}\\[4pt]
\includegraphics[width=0.92\linewidth]{fig/mismatch_robustness_walker.png}
\caption{\textbf{Un-tunable robustness to traction loss.} Planner uses the nominal model;
reality is slippery; no controller is retuned. Top: Hopper. Bottom: Walker2d. A single fixed
\fpl{} spec (red) keeps the highest productive speed and holds survival across a $3.3\times$
friction reduction; the linear family cannot---aggressive weights lose survival, the safe
weight loses speed, and which weight is ``best'' shifts with the friction you don't know.}
\label{fig:robust}
\end{figure}

\begin{figure}[t]
\centering
\includegraphics[width=0.92\linewidth]{fig/quadruped_robustness.png}
\caption{\textbf{Robustness generalizes to a 12-DoF quadruped (Barkour).} Same traction-loss
protocol, no retuning. \fpl{} (red) holds 100\% survival at top-or-tied productive speed across
the sweep; the aggressive linear weights collapse and the near-competitor's survival wobbles
(unpredictable falls).}
\label{fig:robust_quad}
\end{figure}

\paragraph{4. Robustness on a second mismatch axis and off legged robots (Fig.~\ref{fig:robust_cube}).}
The same un-tunable robustness holds on the dexterous hand, with the aggression-punishing
mismatch being \textbf{reduced contact grip} (a worn / slippery cube---the manipulation analogue
of traction loss) instead of ground friction. The hand must still tumble the cube $46\degree$
without dropping it; no controller is retuned across grip levels. One fixed \fpl{} spec keeps
$\sim$100\% survival and the highest productive rotation across a $30\%$ grip reduction, and is
highest at \emph{every} level, its margin \emph{growing} as grip is lost:

\begin{center}\small
\begin{tabular}{lcc}
\toprule
grip (friction) & \fpl{} productive rotation (\degree $\times$ surv) & best fixed linear (wa$=2$) \\
\midrule
1.0 (nominal) & \textbf{26.1 / 100\%} & 18.6 / 100\% \\
0.8           & \textbf{22.6 / 100\%} & 12.9 / \phantom{0}77\% \\
0.7           & \textbf{20.1 / \phantom{0}97\%} & \phantom{0}7.0 / \phantom{0}47\% \\
0.6           & \textbf{12.0 / \phantom{0}67\%} & \phantom{0}5.5 / \phantom{0}30\% \\
\bottomrule
\end{tabular}\\[2pt]
{\footnotesize Productive rotation (\degree{} toward goal $\times$ survival) and survival, 30
seeds. The best fixed linear weight both \emph{loses the cube} (survival $100\%\to30\%$) and sees
its rotation collapse; aggressive weights are dead by grip $0.9$. This is a \emph{second}
aggression-punishing mismatch (grip, not traction) on a \emph{fourth}, non-locomotion robot.}
\end{center}

\begin{figure}[t]
\centering
\includegraphics[width=0.92\linewidth]{fig/cube_robustness.png}
\caption{\textbf{Robustness generalizes off legged robots (LEAP hand, grip loss).} In-hand cube
reorientation; the true contact friction is reduced (slippery cube), no retuning. \fpl{} (red)
holds $\sim$100\% survival and the highest productive rotation across a $30\%$ grip reduction; the
whole linear family collapses (best weight's survival $100\%\to30\%$). Four robots and two
mismatch axes (ground traction + contact grip), same phenomenon.}
\label{fig:robust_cube}
\end{figure}

\paragraph{5. The advantage is the \emph{objective}, not the sampler (Fig.~\ref{fig:portability}).}
The obvious objection to any cost-side result is that a smarter \emph{sampler} with a plain linear
cost might close the gap. We rule this out on both sides. \textbf{(a) Portability:} holding the
sampler fixed and varying only the composition, one fixed \fpl{} spec Pareto-dominates the
linear-weight frontier \emph{within every sampler we tried}---plain Gaussian \mppi, \mppi-CMA,
cross-entropy (CEM), and a colored-noise sampler (16 seeds; walker $+17$--$34\%$, hopper
$+29$--$209\%$ over the fastest comparably-safe linear weight; on the contact-rich cube \fpl{} wins
for \mppi/CEM/colored and \emph{ties} \mppi-CMA, the one honest exception). The un-tunable traction
robustness likewise holds \emph{inside} every sampler (Fig.~\ref{fig:portability}, bottom).
\textbf{(b) No sampler beats plain \mppi.} Reciprocally, holding the \fpl{} cost fixed and varying
only the proposal---including a new FPL-native colored-noise + absolute-scale sampler and
gradient-guided (BPTT) refinement---\emph{nothing beats warm-started plain Gaussian \mppi} at any
budget (Fig.~\ref{fig:samplerrace}; 16 seeds $\times$ 5 budgets). So the entire advantage lives in
the \emph{objective algebra}, and it transfers to whatever sampler you already use---which is the
cleanest possible statement of ``change the scalarization, not the search.''

\begin{figure}[t]
\centering
\includegraphics[width=\linewidth]{fig/fpl_portability.png}\\[4pt]
\includegraphics[width=\linewidth]{fig/fpl_robustness_samplers.png}
\caption{\textbf{The \fpl{} win is the objective, and it is portable across samplers.} Top: within
each of \mppi, \mppi-CMA, CEM, and a colored-noise sampler, the single fixed \fpl{} spec (red star)
sits above-and-right of the whole linear-weight frontier (walker, hopper). Bottom: the un-tunable
traction-loss robustness holds inside every sampler too---\fpl{} rides above the linear family
across the friction sweep in all four panels (hopper, no retuning, 12 seeds, 95\% CIs).}
\label{fig:portability}
\end{figure}

\begin{figure}[t]
\centering
\includegraphics[width=\linewidth]{fig/fpl_sampler_race.png}
\caption{\textbf{No sampler beats plain \mppi{} (a clean negative that anchors the message).} With
the \fpl{} cost held fixed, five proposals---plain Gaussian \mppi, \mppi-CMA, CEM, colored noise,
and colored noise $+$ an FPL absolute-scale schedule---are within CIs at every sample budget, and
plain \mppi{} is best-or-tied throughout; gradient-guided refinement (not shown) is $\sim$100$\times$
costlier and no better. Across four search mechanisms nothing beats warm-started Gaussian \mppi.}
\label{fig:samplerrace}
\end{figure}

\paragraph{6. Off-the-shelf across morphologies, no tuning (Fig.~\ref{fig:zerotuning}).}
Because the \fpl{} objective carries \emph{no} cross-objective trade-off weights---each atom is a
bounded fulfillment with a satisfied-band in physical units, and the min-conjunction composes
them---a \emph{single} \fpl{} spec transfers across robots untuned. Dropped on the walker, hopper,
and hand with no retuning, it is the most productive objective on \emph{every} robot; no single
global linear weight transfers (an aggressive weight runs the walker but falls on the hopper at
$0\%$ survival and drops the cube; a timid weight moves nothing). By contrast the linear cost's
weights are arbitrary and entangled (e.g.\ the hopper cost hand-tunes four ratios $10/50/5/0.3$ with
no natural scale, which must be re-tuned per robot and per difficulty). This is the ``drop it on a
new morphology and get an honest read on whether it can do the task, with no hyperparameter hunt''
use case that motivates the method.

\begin{figure}[t]
\centering
\includegraphics[width=0.82\linewidth]{fig/fpl_zero_tuning.png}
\caption{\textbf{Off-the-shelf: one global \fpl{} spec is the most productive objective on every
robot (boxed); no single linear weight transfers.} Cell $=$ progress (\% of \fpl) $\times$
survival; 16 seeds. Aggressive linear weights (wv$=4$/$8$) run the walker but fall on the hopper
($0\%$ survival) and drop the cube; timid weights (wv$\le 1$) are too slow everywhere.}
\label{fig:zerotuning}
\end{figure}

\paragraph{7. We probed the rest of the \mppi{} pipeline; the objective is where the value is.}
Beyond the cost we tried to give \fpl{} an edge at two further pipeline stages, and report both
honestly. \emph{Weighting/temperature:} \fpl's bounded reward lets the importance weights target a
fixed effective-sample-size instead of a hand-tuned temperature; this genuinely helps where a fixed
temperature collapses to near-uniform weights ($+13\%$ productive speed, $100\%$ vs $92\%$ survival
on the noisy hopper), but the effect is modest and regime-dependent. \emph{Constraint/shield:} the
per-objective structure uniquely enables a hard lexicographic safety filter (satisfy safety first,
then optimize performance)---something a scalar weighted sum cannot express---but for these
short-horizon dynamic tasks it does \emph{not} beat the soft min-fulfillment cost: a hard constraint
under a short horizon is boundary-seeking (a recursive-feasibility failure) and the soft
conjunction, which balances continuously, wins. This is an informative negative: it shows the soft
min-fulfillment is the \emph{right} \fpl{} design here, and localizes where a hard shield / terminal
value would pay off---long horizons or a learned value---i.e.\ exactly the teacher/student RL
extension.

\section{Why this is publishable}

\begin{enumerate}[leftmargin=1.4em,itemsep=3pt]
\item \claim{It answers the obvious objection.} A performance edge on a tradeoff curve always
      invites ``retune the baseline.'' The robustness result is defined precisely by the
      \emph{impossibility} of retuning: the mismatch is unknown at design time. This converts
      ``\fpl{} is a better tuning'' into ``\fpl{} needs no tuning and a linear cost
      \emph{cannot} be tuned for this,'' which is a structural claim.
\item \claim{It matches theory the field already believes.} The information-theoretic \mppi{}
      literature (Williams et al.) fixes the sampling distribution and hand-weights the cost;
      robust/tube \mppi{} adds machinery to the \emph{sampler/dynamics}. We instead change the
      \emph{objective algebra} and show it buys robustness for free, which is an unusual and
      clean lever. \fpl{}'s minimum-fulfillment bound predicts exactly the observed behavior
      (a high conjunction guarantees a floor on every objective), giving the empirics a
      mechanism rather than a coincidence.
\item \claim{The negative results are a feature.} We show \emph{when} \fpl{} does not help
      (static tasks, uniform-difficulty mismatch such as added mass) and give a one-line
      predictive rule: \fpl{} wins iff the mismatch/competition makes \emph{aggression
      dangerous while safe progress remains possible}. A method paper that states its own
      boundary is more, not less, credible---and the rule is testable.
\item \claim{Reproducible and cheap.} Pure-CPU \mppi{}, standard MuJoCo robots, self-contained
      per-result scripts, and a \emph{crash-safe resumable} trial store: every
      (experiment, sampler, cost, seed) trial is one durable line, so all $\sim$3{,}900 trials
      behind the figures regenerate from cache and a killed run resumes with no recomputation. No
      learned components to muddy attribution. Every claim is also \textbf{video-verified}:
      synchronized side-by-side clips (\fpl{} vs.\ best-safe linear vs.\ aggressive linear) for all
      four robots at \texttt{runs/paper/\{walker,hopper,cube,quadruped\}\_sidebyside.mp4}.
\end{enumerate}

\section{Where it is weak (so we decide with eyes open)}

\begin{itemize}[leftmargin=1.4em,itemsep=2pt]
\item \emph{Most tasks are locomotion} (Hopper, Walker2d, a 12-DoF quadruped), \emph{but no longer
      all}: the 16-DoF LEAP-hand in-hand cube reorientation is a non-locomotion
      dynamic-competition + robustness win, so the phenomenon is not specific to legged gaits.
      The remaining gap is a humanoid we could not get to walk from scratch with this planner, so
      that flashiest demo is absent.
\item \emph{Scope of the robustness win is bounded to aggression-punishing mismatch}: it holds
      for two such axes now (ground traction \emph{and} contact grip), but not for
      uniform-difficulty mismatch (mass, actuator strength). We must state this plainly;
      over-claiming ``\fpl{} is universally robust'' would be false.
\item \emph{No pipeline stage beyond the objective helped much (which we now frame as a strength).}
      We tried, and honestly report, FPL-native mechanisms at the sampler (colored noise,
      absolute-scale schedule, adaptive covariance, gradient refinement), the weighting/temperature
      (adaptive effective-sample-size, feasibility gate), and the constraint stage (a hard
      lexicographic safety shield). The sampler and shield are clean negatives and the weighting is
      only a modest, regime-dependent gain (\S5--7). This trims sections we hoped to write, but
      tightens the message to \emph{scalarization, not search}.
\item \emph{Two honest caveats in the new sampler-axis results.} On the contact-rich cube, the
      portability win holds for \mppi/CEM/colored but \emph{ties} \mppi-CMA (its covariance
      adaptation lets even the conservative linear weight rotate safely); and the quadruped is
      absent from the sampler/zero-tuning matrices because its $343$\,s/episode makes a
      multi-sampler grid impractical---its objective-level win stands at 30 seeds
      (Figs.~\ref{fig:pareto},~\ref{fig:robust_quad}).
\item \emph{Perfect-model planning otherwise.} Apart from the mismatch study we plan on the true
      model; the method matters least there, which is why the mismatch study must be the anchor.
\end{itemize}

\section{Positioning and venue}

Framed on \emph{performance}, this is a borderline workshop paper (a nice tradeoff-curve
result). Framed on \emph{un-tunable robustness under model mismatch}, with the honest boundary
and the mechanism, it is a solid short conference / workshop paper in the learning-for-control
space (e.g.\ L4DC, CoRL workshop, ICRA/IROS). The title should foreground robustness, e.g.\
``\emph{Objective algebra, not weight tuning: fulfillment-priority costs give retuning-free
robustness in sampling-based MPC.}'' The teacher/student distillation angle (one objective
spec shared by an \fpl{} \mppi{} teacher and an RL student) is the natural full-paper extension
but is out of scope for a first submission.

\section{What would make it bullet-proof (a to-do for reviewers)}

\begin{enumerate}[leftmargin=1.4em,itemsep=2pt]
\item \sout{A third robot} \sout{\textbf{[done: a 12-DoF quadruped]}} \sout{Next: a
      \emph{non-locomotion} dynamic-competition task} \textbf{[both done: a 12-DoF quadruped and a
      16-DoF LEAP hand (in-hand cube reorientation) each show the dynamic-competition dominance
      \emph{and} the robustness; the hand is non-locomotion, closing the ``legged gaits only''
      gap]}.
\item The minimum-fulfillment bound stated as a proposition and checked numerically against the
      observed per-objective floors, tying theory to the survival curves.
\item \sout{A second mismatch axis that \emph{also} punishes aggression} \textbf{[partly done:
      contact-grip loss on the hand is a second aggression-punishing axis beyond ground traction;
      a sudden external push / disturbance-rejection axis is still worth adding]}.
\item \sout{Seeds bumped to $\ge 30$ and confidence intervals on every point}
      \textbf{[done: all sweeps are 30 seeds with 95\% CIs — x/y error bars on the Pareto
      plots, shaded bands on the robustness plots (Wilson intervals for survival)]}.
\end{enumerate}

\paragraph{My recommendation.} Write it, lead with Fig.~\ref{fig:robust}, keep the nulls in
the main text. With item (1) done (four robots incl.\ a non-locomotion hand, two mismatch axes) and
the new \emph{objective-not-search} evidence (portability within every sampler,
Fig.~\ref{fig:portability}; no sampler beats plain \mppi, Fig.~\ref{fig:samplerrace}; and
off-the-shelf transfer with no tuning, Fig.~\ref{fig:zerotuning}), the empirical case is now
comfortably conference-ready and arguably a short full paper. The remaining polish is item (2) (the
minimum-fulfillment bound as a checked proposition) and a disturbance-rejection axis. The robustness
and objective-portability results are the real contribution; the honest pipeline-stage negatives
(\S7) are what make it believable; the teacher/student RL extension is the natural full-paper
sequel that the shield negative directly motivates.

\appendix
\section{Supplementary material (all regenerable / checkpointed)}
\textbf{Companion write-up:} \texttt{FPL\_MPPI\_CAMPAIGN.md} (this campaign's full results) and
\texttt{FPL\_FINDINGS.md} (the objective-level wins). \textbf{Additional figures}
(\texttt{fig/}): \texttt{fpl\_portability\_cube.png} (cube portability per sampler at
$K{=}256$), \texttt{fpl\_weighting.png} $+$ \texttt{fpl\_weighting\_robustness.png} (the weighting
layer, \S7), \texttt{fpl\_shield.png} (the hard-safety shield negative, \S7).
\textbf{Videos} (\texttt{runs/paper/}): synchronized side-by-side
\texttt{\{walker,hopper,cube,quadruped\}\_sidebyside.mp4} and single-controller
\texttt{\{...\}\_\{fpl,linear\}.mp4}. \textbf{New controllers} (\texttt{analytic\_mppi/controllers/}):
\texttt{fpl\_colored} (colored-noise sampler), \texttt{fpl\_tempered} (adaptive-ESS weighting),
\texttt{fpl\_shielded} (lexicographic safety filter). \textbf{Scripts} (\texttt{verification/}):
\texttt{fpl\_sampler\_race.py}, \texttt{fpl\_portability.py}, \texttt{fpl\_zero\_tuning.py},
\texttt{fpl\_robustness\_samplers.py}, \texttt{fpl\_weighting\_study.py},
\texttt{fpl\_weighting\_robustness.py}, \texttt{fpl\_shield\_study.py},
\texttt{fpl\_gradient\_probe.py}, \texttt{g1\_reach\_aggressive.py}; all write resumable
\texttt{verification/checkpoints/*.jsonl}.

\end{document}
